\chapter{Design, Testing and Evaluation}
\label{ch:design}

\section{Server Design}

The server is split into three files, and each one has a single job:
\texttt{db.js} handles the database and the security helper functions;
\texttt{server.js} defines the web app and every API route; and
\texttt{websocket.js} handles the live connection and the function that
tells a waiting laptop that its login was approved. Splitting the code
this way keeps each file easy to read on its own, rather than having one
huge file that does everything.

\subsection{Database Layout}

Table~\ref{tab:schema} lists the seven tables that make up Echo's
database. Two of them, \texttt{login\_\allowbreak sessions} and \texttt{enroll\_\allowbreak tokens},
are cleaned up the most carefully: both have an expiry time and a used
flag, and a background task (Section~\ref{sec:architecture}) deletes
expired, unused rows every 60 seconds so the database does not fill up
with old, useless data over time.

\begin{table}[htbp]
\centering
\caption{The database tables}
\label{tab:schema}
\begin{tabular}{|L{3cm}|L{9.5cm}|}
\hline
\textbf{Table} & \textbf{What it stores} \\
\hline
\texttt{users} & Username, an optional password hash (kept for older accounts), optional email, when the account was created \\
\hline
\texttt{devices} & One row per registered phone: which user owns it, its name, and its public key \\
\hline
\texttt{login\_sessions} & One row per login attempt: the one-time code, its status, a one-time claim token, when it expires, whether it has been used \\
\hline
\texttt{sessions} & Logged-in session tokens (the cookie value), with an expiry time \\
\hline
\texttt{enroll\_tokens} & One-time sign-up codes, tied to a username, with an expiry time and a used flag \\
\hline
\texttt{recovery\_codes} & Scrambled (hashed) versions of one-time recovery codes, from the older recovery method \\
\hline
\texttt{logins} & A running log of every login attempt: user, method used, phone, whether it worked, and when \\
\hline
\texttt{magic\_tokens} & One-time email recovery links, with an expiry time and a used flag \\
\hline
\end{tabular}
\end{table}

Where a password hash is stored at all, it is scrambled using a standard,
slow method called PBKDF2, with a random 256-bit salt and 310{,}000
rounds, following the current minimum recommended by NIST (the body that
sets US security standards) for this method. Checking a password always
runs the full scrambling process and compares the result carefully, rather
than stopping early, so that the time it takes cannot leak any hints about
whether a guess was close. Recovery codes and email links are likewise
never stored as plain, readable text.

\subsection{The API}

The server offers a set of web addresses (an API) under \texttt{/api},
plus one live WebSocket connection. Table~\ref{tab:api} lists them all;
the full details of what each one expects and returns are in
Appendix~\ref{app:api}.

\begin{table}[htbp]
\centering
\caption{The API, in brief}
\label{tab:api}
\begin{tabular}{|L{4.9cm}|L{7.6cm}|}
\hline
\textbf{Address} & \textbf{What it does} \\
\hline
\texttt{POST /api/\allowbreak signup} & Claim a username, get back a one-time sign-up code \\
\hline
\texttt{POST /api/\allowbreak enroll} & Use that code to register a phone's public key \\
\hline
\texttt{GET /api/\allowbreak signup/\allowbreak status} & Check whether sign-up has finished yet \\
\hline
\texttt{POST /api/\allowbreak login/\allowbreak start} & Create a one-time login code to play as sound \\
\hline
\texttt{GET /api/\allowbreak login/\allowbreak check} & The phone's quiet check that a code belongs to its own account \\
\hline
\texttt{POST /api/\allowbreak login/\allowbreak verify} & Send the phone's signature; approves the login if it checks out \\
\hline
\texttt{POST /api/\allowbreak session/\allowbreak claim} & Swap a one-time claim token for a real login cookie \\
\hline
\texttt{POST /api/\allowbreak login/\allowbreak magic-request} & Ask for an email recovery link \\
\hline
\texttt{GET /api/\allowbreak login/\allowbreak magic} & Use that email link to log in \\
\hline
\texttt{POST /api/\allowbreak device/\allowbreak token} & Get a code to register a second phone (must be logged in) \\
\hline
\texttt{POST /api/\allowbreak device/\allowbreak revoke} & Remove a registered phone (must be logged in) \\
\hline
\texttt{GET /api/\allowbreak me} & See your profile, phones, and recent logins (must be logged in) \\
\hline
\texttt{POST /api/\allowbreak logout} & End the current login session \\
\hline
\texttt{WS /ws?session=<id>} & The laptop's live connection, used to hear "you're approved" \\
\hline
\end{tabular}
\end{table}

\subsection{The Laptop Page and the Phone App}

The laptop login page is a simple web page that asks the server for a
one-time code, opens a live connection for that login attempt, and uses
\texttt{ggwave} to turn the code into sound and play it out loud (though
too high-pitched to actually hear). A status message on screen moves
through three steps, sending, phone heard it, approved, so the person
watching can tell what stage the login has reached, rather than staring
at a blank screen for up to 30 seconds. If nothing happens, the page tries
again automatically, up to three times.

The phone app installs itself from the browser onto the phone's home
screen, and can even work offline after the first time it is opened. Once
open, it asks for permission to use the microphone and listens
constantly for a sound signal. When it picks one up, it runs the quiet
background check described in Section~\ref{sec:methodology-crypto} before
showing anything on screen, so the phone only ever asks its owner to
approve logins for accounts it is actually registered to. Tapping approve
makes the phone sign the login using its private key, which is stored
safely on the phone and never leaves it, and send that signature straight
to the server.

\subsection{How the Design Changed During the Project}
\label{sec:design-evolution}

Because Echo was built step by step (Chapter~\ref{ch:methodology}), a few
details changed from the original plan once the working system showed
real problems the plan on paper had missed. Looking back at the project
history against the original plan shows three clear changes, listed here
honestly rather than left out:

\begin{itemize}
  \item \textbf{How account recovery works.} The original plan called for
  six one-time numeric recovery codes per account, stored as scrambled
  hashes. That was built (the \texttt{recovery\_codes} table and its code
  are still in the project), but it was replaced as the main recovery
  method by a rate-limited email link instead, which felt more practical
  for someone without a pre-printed list of codes on hand. The key safety
  property stayed the same either way: recovery is rate-limited and cannot
  be brute-forced.
  \item \textbf{How sign-up works.} The original plan assumed the sign-up
  code would be typed into the phone by hand. A QR-code scanner was added
  to the phone app instead, so the laptop can show the sign-up link as a
  scannable code. This made signing up noticeably faster without changing
  how the underlying code actually works or how safe it is.
  \item \textbf{Making the live demo more reliable.} Hosting the project
  on Render.com's free plan brought up a real-world problem the original
  plan did not consider: the server falls asleep after about 15 minutes
  with no visitors, so the very first request after that would take up to
  30 seconds to load. A small background task was added that pings the
  live server every few minutes to stop it from falling asleep during
  demonstrations. It has no effect when running the project locally.
\end{itemize}

Other work, such as a dark mode toggle, a more polished home page, and
consistent styling across the site, was also done during the project but
does not change how the login itself works or how secure it is, so it is
not covered in detail here.

\section{Testing}
\label{sec:testing}

The correctness of the login process was checked using automated tests
that talk to a running copy of the server over the network, using real
signing keys to both sign and forge messages, but no actual sound
hardware, so the whole set of tests runs in seconds and can be repeated
after every change. The tests are split into two files:
\texttt{test-flow.js} checks the whole login process from start to
finish, including the attack scenarios that matter most for the security
claims in Section~\ref{sec:methodology-crypto}, and
\texttt{test-cross-talk.js} checks specifically that one user's phone can
never be asked to approve a different user's login.

Both files were run again against the current code for this report, using
Node.js version 22.22.3, on 8 August 2026. Table~\ref{tab:testresults}
shows the results. Every check that matters for the security claims in
this report, replaying a recording, forging a signature, one user's device
being used against another account, expired codes, and reused tokens, all
passed.

\begin{table}[htbp]
\centering
\caption{Automated test results (Node.js v22.22.3, run locally)}
\label{tab:testresults}
\begin{tabular}{|L{6.3cm}|L{2.3cm}|L{2.3cm}|}
\hline
\textbf{Test file} & \textbf{Checks} & \textbf{Result} \\
\hline
\texttt{tests/test-flow.js} (full login process, including replay, forgery, wrong-device, reused claim token, recovery, removing a phone) & 22 & 22 passed, 0 failed \\
\hline
\texttt{tests/test-cross-talk.js} (stops one user's phone approving another user's login) & 9 & 9 passed, 0 failed \\
\hline
\textbf{Total} & \textbf{31} & \textbf{31 passed, 0 failed} \\
\hline
\end{tabular}
\end{table}

Looking specifically at the threat model in
Section~\ref{sec:methodology-crypto}, these attack attempts were tested
directly and every single one was correctly blocked: signing up without a
valid code; reusing an already-used sign-up code; replaying an
already-used login code; a signature made with the wrong private key; a
correct signature used against the wrong user's device; reusing a claim
token that has already been used once; and trying to log in with a phone
that has been removed from the account. Two smaller checks also confirm
that only a phone's own account passes its quiet background check, and
that a made-up code and an already-used code are correctly told apart
with different, clear error messages, which matters for spotting problems
quickly during a live demo.

\section{Evaluation}
\label{sec:evaluation}

\subsection{How the Real-World Test Was Planned}

Section~\ref{sec:methodology-eval} set a target of at least 90\% success
at 1\,m in a quiet room, to be tested across nine combinations: three
distances (0.3\,m, 1\,m, 2\,m) and three noise levels (quiet, speech,
music), with at least 30 attempts for each combination, 270 attempts in
total. Each attempt records: whether the phone picked up the code at all,
whether the whole login finished successfully, how long it took from
start to finish, the distance, and the noise level, matching the format
already planned for in the project's evaluation tool
(Appendix~\ref{app:api}).

\begin{table}[htbp]
\centering
\caption{Real-world test results, to be filled in (270 attempts: 3 distances times 3 noise levels times 30 attempts each). This table should be filled in with real measurements from the actual laptop and phone before the report is submitted.}
\label{tab:evalresults}
\begin{tabular}{|L{2.6cm}|L{2.4cm}|L{2.2cm}|L{2.4cm}|L{2.4cm}|}
\hline
\textbf{Distance} & \textbf{Noise} & \textbf{Attempts} & \textbf{Success rate} & \textbf{Typical time} \\
\hline
0.3 m & Quiet & 30 & \emph{to be filled in} & \emph{to be filled in} \\
\hline
0.3 m & Speech & 30 & \emph{to be filled in} & \emph{to be filled in} \\
\hline
0.3 m & Music & 30 & \emph{to be filled in} & \emph{to be filled in} \\
\hline
1 m & Quiet & 30 & \emph{to be filled in} & \emph{to be filled in} \\
\hline
1 m & Speech & 30 & \emph{to be filled in} & \emph{to be filled in} \\
\hline
1 m & Music & 30 & \emph{to be filled in} & \emph{to be filled in} \\
\hline
2 m & Quiet & 30 & \emph{to be filled in} & \emph{to be filled in} \\
\hline
2 m & Speech & 30 & \emph{to be filled in} & \emph{to be filled in} \\
\hline
2 m & Music & 30 & \emph{to be filled in} & \emph{to be filled in} \\
\hline
\end{tabular}
\end{table}

\textbf{A note on this table.} Table~\ref{tab:evalresults} is left as a
template on purpose, not filled in with made-up numbers. How well the
sound carries, and how long a login takes, depends on the exact laptop
speaker and phone microphone used, the shape and material of the room, and
the background noise on the day of testing. None of that can be measured
while writing this report on a computer with no speakers or microphone to
test with. The plan for collecting this data (recording success, timing,
distance, and noise for every attempt) is ready to use; what is left is to
actually run the 270 test attempts on the real laptop and phone, fill in
this table with the real numbers, and update the conclusion in
Section~\ref{sec:conclusion-findings} to match. This gap is stated
plainly here, instead of being hidden behind invented numbers, because
made-up results would give a false picture of how well the system
actually performs in the real world, which is the whole point of this
section.

\subsection{An Early Sign It Works}

Separate from the full 270-attempt test, an early feasibility test
(\texttt{tests/\allowbreak ultrasonic-\allowbreak auth-\allowbreak test.html}, mentioned in
Chapter~\ref{ch:methodology}) already showed that the 18 to 20\,kHz sound
channel works between a real laptop speaker and a real phone microphone
at close range. That result is the reason the project moved ahead into
full development in the first place. That same test page is still
available for quick, informal checks, and it was used to help choose the
0.3 to 2\,m distance range and the quiet/speech/music noise levels used in
the full test plan above.
